\scanpage{71}

% top center and top right
\mspage{- 140 -}\hfill IX

% upper, definition and first relation
où par définition $({}_nA_K)^f=({}_nA^{\unsure{\mathrm{o}|0}})^f_K$ désigne le r.ème composant $(n=\ell^{r+1})$ du pro-groupe de Barsotti-Tate $T_\ell(A_K)^f$. La première des relations (5.5.8) nous fournit alors un isomorphisme canonique

% upper, (11.6.7)
\[
{}_n\Psi_K\simeq\underline{M}_{nK}\qquad\text{où}\quad\underline{M}_n\overset{\mathrm{dfn}}{=}\underline{M}\otimes_{\mathbb{Z}}(\mathbb{Z}/n\mathbb{Z})\ .\tag*{(11.6.7)}
\]

% upper middle, transition
En fait, on a :

% upper middle, lemma heading and assertion
\uline{Lemme} 11.6.8. \uline{Le schéma en groupes} $ {}_n\Psi$ (11.6.6) \uline{sur} $S$ \uline{est étale sur} $S$, \uline{et l'isomorphisme} (11.6.7) \uline{se prolonge de façon unique en un morphisme}

% middle, open immersion
\[
{}_n\Psi\handwritten{\hookrightarrow}\underline{M}_n\tag*{(11.6.8.1)}
\]

% middle, lemma's last line
\uline{qui est une immersion ouverte}.

% middle and lower, proof of lemma
Pour voir que $ {}_n\Psi$ est étale, il reste à vérifier en vertu de (11.6.7) qu'il l'est en les points de sa fibre spéciale, ou encore que cette dernière est étale, ce qui est trivial, la fibre en question étant $ {}_nA_{\mathrm{o}}/{}_nA_{\mathrm{o}}^{\mathrm{o}}\subset A_{\mathrm{o}}/A_{\mathrm{o}}^{\mathrm{o}}=\Phi_{\mathrm{o}}$. Il s'ensuit que $ {}_n\Psi$ est un schéma normal, et comme $\underline{M}_n$ est fini sur $S$, l'existence d'un prolongement (11.6.8.1) en résulte aussitôt. C'est un morphisme de schémas en groupes étales sur $S$, donc son noyau est un groupe étale sur $S$, et comme $ {}_n\Psi$ donc ce noyau est séparé, et sa fibre générique est le groupe unité par (11.6.7), c'est le groupe unité (3.1.3). Donc (11.6.8.1) est un monomorphisme étale, donc une immersion ouverte (EGA IV 17.9.1), ce qui prouve le lemme.

% lower, new paragraph
Posons maintenant

% lower, definition (11.6.9)
\[
{}_n\Phi=({}_nA)^f/({}_nA^{\mathrm{o}})^f\subset{}_n\Psi\quad,\tag*{(11.6.9)}
\]

% lower, continued sentence
de sorte que $ {}_n\Phi$ est un sous-schéma fini ouvert de $ {}_n\Psi$, puisque $({}_nA)^f$

% bottom strip, pencil numeral
\handwritten{458}

% bottom center, printed number on label
458

% lower left
% archive sheet number: 70
